% =====================================================================
%  explainability.tex  --  PHASE VI, Tasks 13A-13C
%
%  Shapley attribution of the two forecasting components and of the
%  state layer, and the transition structure the temporal model learns.
%
%  Intended placement: Section 7 (\label{sec:explain}), after the
%  decision layer and before the ablation.
%
%  Drop-in usage: % =====================================================================
%  explainability.tex  --  PHASE VI, Tasks 13A-13C
%
%  Shapley attribution of the two forecasting components and of the
%  state layer, and the transition structure the temporal model learns.
%
%  Intended placement: Section 7 (\label{sec:explain}), after the
%  decision layer and before the ablation.
%
%  Drop-in usage: % =====================================================================
%  explainability.tex  --  PHASE VI, Tasks 13A-13C
%
%  Shapley attribution of the two forecasting components and of the
%  state layer, and the transition structure the temporal model learns.
%
%  Intended placement: Section 7 (\label{sec:explain}), after the
%  decision layer and before the ablation.
%
%  Drop-in usage: % =====================================================================
%  explainability.tex  --  PHASE VI, Tasks 13A-13C
%
%  Shapley attribution of the two forecasting components and of the
%  state layer, and the transition structure the temporal model learns.
%
%  Intended placement: Section 7 (\label{sec:explain}), after the
%  decision layer and before the ablation.
%
%  Drop-in usage: \input{explainability}
%  Figures are referenced by bare name; add to the preamble
%      \graphicspath{{../outputs/phase6/figures/}}
%  and use the PDF copies (both PDF and PNG are written).
%  Bibliography: paper/references.bib
%  Numbers below are generated by  python run_phase6.py  and stored in
%    outputs/phase6/task13a/<machine>/task13a_results.json
%    outputs/phase6/task13b/<machine>/task13b_results.json
%    outputs/phase6/task13c/task13c_results.json
%  Figures:
%    fig_task13a_decision_shap_phase4, fig_task13a_window_shap,
%    fig_task13b_state_shap_phase4, fig_task13c_transition_matrix,
%    fig_task13c_transition_graph
%
%  FORWARD REFERENCES.  This file cites four sections that do not yet
%  have a .tex of their own; they will resolve to ?? until those are
%  written, which is deliberate rather than an oversight.  The labels
%  expected are:
%    sec:forecasting    the Task 5 / Task 12A forecasting comparison
%    sec:significance   the Phase V significance testing
%    sec:multimachine   the Task 8 eight-machine analysis
%    sec:crossmachine   the Task 9 transfer study
%  Everything else resolves against method_gmm_labeller.tex,
%  related_work.tex and decision_layer.tex as they stand.
% =====================================================================

\section{What the Models Use}
\label{sec:explain}

The three stages of the framework are opaque in different ways. The state
layer assigns a label without saying which measured quantity carried the
decision; both forecasters report split-gain importances, which are known to
be biased toward high-cardinality splits and describe the fitted trees rather
than any individual prediction. This section replaces those importances with
Shapley attributions~\cite{lundberg2017shap}, which are additive, consistent,
and defined per prediction, and reports the transition structure the temporal
model learns.

One caveat governs everything below and is not repeated at each result. Every
target attributed here is derived from the pipeline's own GMM--HMM labels.
These attributions therefore explain what each model learned; they are not
evidence that the underlying state assignment is correct. Validation of the
state definition is the physics battery of
Section~\ref{sec:method:accept_reject} and the cross-machine evidence of
Section~\ref{sec:crossmachine}, not this section.

\subsection{Attribution method}
\label{sec:explain:method}

Every model attributed here is a gradient-boosted
tree~\cite{chen2016xgboost} --- the two forecasters, the chance-constraint
classifier, and the surrogate fitted to the state labels --- so exact tree SHAP
applies and no sampling approximation is needed. Attributions are
computed on held-out rows: $7\,440$ decision epochs for the duration
forecaster, $15\,000$ design rows for the window forecaster, and $20\,000$
energised rows for the state surrogate.

Shapley values satisfy $f(x)=\mathbb{E}[f]+\sum_j\phi_j(x)$, and we verify that
identity rather than assume it. The verification is performed on the model's
raw margin --- the quantity tree SHAP decomposes --- and on a \emph{relative}
scale: the duration forecaster's outputs are seconds and run to $10^{5}$,
while the tree traversal accumulates in single precision, so an absolute
tolerance written for probability outputs is not the right instrument. The
largest relative reconstruction error measured over the four forecaster
attributions is $2.0\times10^{-6}$, four orders of magnitude below the
smallest difference any claim below rests on; the state surrogate's
attributions satisfy the stricter absolute check unmodified.

\subsection{The decision layer's forecaster is a calendar model}
\label{sec:explain:decision}

Section~\ref{sec:crossmachine} reported that the remaining-duration forecaster
draws $94.5\%$ of its split gain from four calendar features, and inferred
from that a model carrying the plant timetable rather than machine physics.
The Shapley attribution confirms the conclusion and corrects the number.

Under the minimum-dwell labelling the seven calendar features carry $82.4\%$
of the mean $|\phi|$ against $96.7\%$ of the split gain; on the
Phase~I--III labelling the same quantities are $78.0\%$ and $96.2\%$. The
direction is the same on both labellings, and the gap between the two measures
is systematic rather than incidental: gain and Shapley agree only at
Spearman $\rho=0.72$ over the ten features ($\rho=0.43$ on the Phase~I--III
labelling), and the two that move furthest are the two the original claim
leaned on.
Figure~\ref{fig:shap-decision}B makes the disagreement visible. The binary
\texttt{cur\_is\_weekend} takes $63.2\%$ of the split gain and $5.1\%$ of the
attribution: it is split on early and often, near the root, on a partition that
subsequent splits then refine, so it accumulates gain out of proportion to the
change it makes to any prediction. Conversely \texttt{onset\_hour} carries
$17.5\%$ of the gain and $43.1\%$ of the attribution.

\textbf{The corrected statement of the finding is therefore that four fifths,
not nineteen twentieths, of the decision forecaster's behaviour is calendar
--- and the conclusion drawn from it in Section~\ref{sec:crossmachine} is
unaffected.} A model that is $82\%$ plant timetable still transfers between any
two machines sharing a shift pattern and still transfers to nothing else,
which is exactly what the transfer matrices measure. What the attribution adds
is that the machine's own history is not negligible: the previous production
run contributes $13.3\%$ of the attribution against $2.3\%$ of the gain, so the
forecaster is not purely a clock.

The chance-constraint classifier keys on something different, and the
difference is what justifies fitting it separately rather than deriving it
from the point forecast. Elapsed idle time is its single largest input
($33.1\%$ of attribution, rank 1) against sixth and $4.3\%$ for the duration
regressor. The feasibility question --- \emph{will this episode last at least
$h_{\min}$ longer?} --- depends on how long the episode has already run in a
way the conditional mean does not.

\begin{figure}[t]
  \centering
  \includegraphics[width=\linewidth]{fig_task13a_decision_shap_phase4}
  \caption{Shapley attribution of the remaining-idle-duration forecaster on
  the minimum-dwell labelling. (A)~Per-prediction attributions on $7\,440$
  held-out decision epochs; the horizontal axis is symmetric-log because
  attributions span $\pm10^{5}$~s while the bulk lies inside $\pm10^{3}$~s.
  Colour is the feature's own value. (B)~Mean $|\phi|$ against split gain; the
  dashed line is where the two measures would agree.}
  \label{fig:shap-decision}
\end{figure}

\subsection{The window forecaster is an electrical model, and mostly a
persistence one}
\label{sec:explain:window}

The window forecaster and the duration forecaster are separate models serving
separate consumers, and their attributions have almost nothing in common. The
window model draws $40.8\%$ of its attribution from the raw electrical
channels, $23.9\%$ from the rolling statistics and $20.1\%$ from the derived
electrical features, against $11.6\%$ for everything calendar
(Figure~\ref{fig:shap-window}A). Set beside the $82\%$ of the previous
subsection, this is the quantitative form of the observation in
Section~\ref{sec:crossmachine} that the two layers depend on nearly disjoint
information --- and it explains why one transfers between machines and the
other does not.

Two results here are worth stating plainly because neither is convenient.

\textbf{Voltage is the single largest contributor}, at $20.0\%$ of the
attribution across its nine design columns and $11.8\%$ on the final observed
value alone, against $3.0\%$ of the split gain. Supply voltage is not a state
variable of the machine; it sags under load and recovers when the drive
unloads, so it is a legitimate but indirect indicator, and one that would not
survive a change of supply transformer or a different point of common
coupling. A deployment on another site should expect this contribution not to
transfer.

\textbf{Almost half the attribution sits on the final observed value.}
Grouping the design columns by what they summarise
(Figure~\ref{fig:shap-window}B) gives $43.5\%$ to the last sample of the
lookback window, $26.4\%$ to block means and $26.3\%$ to block standard
deviations. A model whose largest single input is the state the machine is in
right now is a refinement of ``nothing changes'', and it cannot be expected to
beat that reference by much. It does not: the measured margin over untrained
persistence is $1.45$~s of per-window MAE
(Section~\ref{sec:significance}), which the attribution here now explains
rather than merely reports.

\begin{figure}[t]
  \centering
  \includegraphics[width=\linewidth]{fig_task13a_window_shap}
  \caption{Shapley attribution of the window state forecaster, refitted at
  seed~0 under the protocol of Section~\ref{sec:forecasting} and verified to
  reproduce its stored metrics exactly. (A)~The twelve largest measured
  channels, summed over their nine design columns each; hollow bars are the
  split gain. (B)~The same attribution grouped by what each design column
  summarises.}
  \label{fig:shap-window}
\end{figure}

\subsection{The state layer separates STANDBY on power factor}
\label{sec:explain:state}

Methods~1 and 2 of the five-method proof rest on a physical assertion: an
energised but unloaded induction motor is identified by its power factor and
its magnetising-current ratio, not by its power magnitude. That assertion is
testable against the labels the pipeline actually produced, by fitting a
gradient-boosted surrogate on them and attributing it.

On the hard split --- STANDBY against $\{$WORKING, PEAK\_LOAD$\}$, restricted
to energised rows so that the trivial OFF boundary cannot dominate --- the
attribution ranks \texttt{current} first ($17.8\%$), \texttt{power\_factor}
second ($16.9\%$) and \texttt{active\_power} third ($12.3\%$). The four
physics-named channels carry $30.9\%$ of the attribution between them.

The result is clearest where the two measures disagree. By split gain,
\texttt{active\_power} takes $60.3\%$ and \texttt{power\_factor} $5.1\%$; by
Shapley attribution the order reverses. \textbf{Power magnitude is what the
trees split on; power factor is what changes the predictions.} That is the
physics argument confirmed by the data rather than asserted from motor theory
--- and it is confirmed on the labelling the paper now quotes, not only on the
one the argument was originally made against. Agreement with the corresponding
Phase~II ranking on the earlier labelling is Spearman $\rho=0.833$
($p=4\times10^{-7}$) with four of five top features shared, so the
minimum-dwell re-labelling does not move the explanation.

Figure~\ref{fig:shap-state}A shows the mechanism directly: for
\texttt{power\_factor} and \texttt{current}, the attributions pushing towards
STANDBY are the low feature values, which is the magnetising-current signature
stated in Section~\ref{sec:method:accept_reject}, read off the model.

The surrogate reproduces its targets to $99.74\%$ accuracy. \emph{That number
is not a validation result and should not be read as one}: the targets are the
state model's own labels, so the accuracy measures only that a tree can
re-express a mixture model's decision boundary. Only the ranking is a result.

\begin{figure}[t]
  \centering
  \includegraphics[width=\linewidth]{fig_task13b_state_shap_phase4}
  \caption{Attribution of a gradient-boosted surrogate fitted on the
  GMM--HMM labels, for the STANDBY-versus-productive split on energised rows.
  (A)~Per-row attributions; blue is a low feature value.
  (B)~Three importance measures with independent failure modes, each
  normalised to its own total.}
  \label{fig:shap-state}
\end{figure}

\subsection{The transition structure}
\label{sec:explain:transitions}

The temporal model is claimed to learn a physically admissible state graph.
Figure~\ref{fig:transitions} reports three matrices that are commonly
conflated: the HMM's estimated transition matrix, the frequencies realised by
Viterbi decoding, and the frequencies after the minimum-dwell constraint.
Transitions are never counted across an acquisition gap.

\textbf{The prediction the state definition makes is that a machine cannot
move between de-energised and loaded in one second.} On pelletizer~I the
learned probability of OFF\,$\rightarrow$\,PEAK\_LOAD is exactly zero, and the
decoded sequence contains \emph{zero} such transitions in $5.47$~million
readings. OFF\,$\rightarrow$\,WORKING occurs $126$ times
($6.1\times10^{-5}$ of the outgoing mass from OFF), which at one-second
sampling is the expected consequence of a machine crossing the
STANDBY band faster than the sampling interval rather than a
violation of the ordering.

\textbf{The minimum-dwell constraint acts on the transitions, not only on the
labels.} WORKING\,$\rightarrow$\,OFF falls from $124$ events to $25$, and the
expected dwell in STANDBY rises from $96$~s to $246$~s. The constraint is
therefore doing what it was introduced to do --- removing runs too short to be
physical --- rather than merely relabelling rows.

It also introduces a small artefact that should be stated: three
PEAK\_LOAD\,$\rightarrow$\,OFF transitions appear after the constraint that
the decoder never produced. Absorbing a sub-threshold intermediate run into
its neighbours can join two states that were not adjacent in the decoded
sequence. Three events in $5.47$~million is $5\times10^{-7}$ of the record and
changes no reported quantity, but it is a property of the operator rather than
of the machine.

Across the plant the structure holds where the state model does and fails where
it does not. The largest forbidden-transition probability among the four
motor-driven machines is $1.9\times10^{-4}$ (pelletizer~I's
WORKING\,$\rightarrow$\,OFF); on exhaust-fan~I it is $3.4\times10^{-3}$,
eighteen times larger --- the same machine whose state names
Section~\ref{sec:multimachine} showed to be mis-assigned, arriving here through
an independent measurement.

\begin{figure}[t]
  \centering
  \includegraphics[width=\linewidth]{fig_task13c_transition_matrix}
  \caption{Transition structure on pelletizer~I. (A)~The HMM's estimated
  matrix. (B)~Frequencies counted from the Viterbi-decoded sequence.
  (C)~The same after the minimum-dwell constraint. Colour is logarithmic;
  parenthesised numbers are event counts for the transitions the state
  definition rules out.}
  \label{fig:transitions}
\end{figure}

%  Figures are referenced by bare name; add to the preamble
%      \graphicspath{{../outputs/phase6/figures/}}
%  and use the PDF copies (both PDF and PNG are written).
%  Bibliography: paper/references.bib
%  Numbers below are generated by  python run_phase6.py  and stored in
%    outputs/phase6/task13a/<machine>/task13a_results.json
%    outputs/phase6/task13b/<machine>/task13b_results.json
%    outputs/phase6/task13c/task13c_results.json
%  Figures:
%    fig_task13a_decision_shap_phase4, fig_task13a_window_shap,
%    fig_task13b_state_shap_phase4, fig_task13c_transition_matrix,
%    fig_task13c_transition_graph
%
%  FORWARD REFERENCES.  This file cites four sections that do not yet
%  have a .tex of their own; they will resolve to ?? until those are
%  written, which is deliberate rather than an oversight.  The labels
%  expected are:
%    sec:forecasting    the Task 5 / Task 12A forecasting comparison
%    sec:significance   the Phase V significance testing
%    sec:multimachine   the Task 8 eight-machine analysis
%    sec:crossmachine   the Task 9 transfer study
%  Everything else resolves against method_gmm_labeller.tex,
%  related_work.tex and decision_layer.tex as they stand.
% =====================================================================

\section{What the Models Use}
\label{sec:explain}

The three stages of the framework are opaque in different ways. The state
layer assigns a label without saying which measured quantity carried the
decision; both forecasters report split-gain importances, which are known to
be biased toward high-cardinality splits and describe the fitted trees rather
than any individual prediction. This section replaces those importances with
Shapley attributions~\cite{lundberg2017shap}, which are additive, consistent,
and defined per prediction, and reports the transition structure the temporal
model learns.

One caveat governs everything below and is not repeated at each result. Every
target attributed here is derived from the pipeline's own GMM--HMM labels.
These attributions therefore explain what each model learned; they are not
evidence that the underlying state assignment is correct. Validation of the
state definition is the physics battery of
Section~\ref{sec:method:accept_reject} and the cross-machine evidence of
Section~\ref{sec:crossmachine}, not this section.

\subsection{Attribution method}
\label{sec:explain:method}

Every model attributed here is a gradient-boosted
tree~\cite{chen2016xgboost} --- the two forecasters, the chance-constraint
classifier, and the surrogate fitted to the state labels --- so exact tree SHAP
applies and no sampling approximation is needed. Attributions are
computed on held-out rows: $7\,440$ decision epochs for the duration
forecaster, $15\,000$ design rows for the window forecaster, and $20\,000$
energised rows for the state surrogate.

Shapley values satisfy $f(x)=\mathbb{E}[f]+\sum_j\phi_j(x)$, and we verify that
identity rather than assume it. The verification is performed on the model's
raw margin --- the quantity tree SHAP decomposes --- and on a \emph{relative}
scale: the duration forecaster's outputs are seconds and run to $10^{5}$,
while the tree traversal accumulates in single precision, so an absolute
tolerance written for probability outputs is not the right instrument. The
largest relative reconstruction error measured over the four forecaster
attributions is $2.0\times10^{-6}$, four orders of magnitude below the
smallest difference any claim below rests on; the state surrogate's
attributions satisfy the stricter absolute check unmodified.

\subsection{The decision layer's forecaster is a calendar model}
\label{sec:explain:decision}

Section~\ref{sec:crossmachine} reported that the remaining-duration forecaster
draws $94.5\%$ of its split gain from four calendar features, and inferred
from that a model carrying the plant timetable rather than machine physics.
The Shapley attribution confirms the conclusion and corrects the number.

Under the minimum-dwell labelling the seven calendar features carry $82.4\%$
of the mean $|\phi|$ against $96.7\%$ of the split gain; on the
Phase~I--III labelling the same quantities are $78.0\%$ and $96.2\%$. The
direction is the same on both labellings, and the gap between the two measures
is systematic rather than incidental: gain and Shapley agree only at
Spearman $\rho=0.72$ over the ten features ($\rho=0.43$ on the Phase~I--III
labelling), and the two that move furthest are the two the original claim
leaned on.
Figure~\ref{fig:shap-decision}B makes the disagreement visible. The binary
\texttt{cur\_is\_weekend} takes $63.2\%$ of the split gain and $5.1\%$ of the
attribution: it is split on early and often, near the root, on a partition that
subsequent splits then refine, so it accumulates gain out of proportion to the
change it makes to any prediction. Conversely \texttt{onset\_hour} carries
$17.5\%$ of the gain and $43.1\%$ of the attribution.

\textbf{The corrected statement of the finding is therefore that four fifths,
not nineteen twentieths, of the decision forecaster's behaviour is calendar
--- and the conclusion drawn from it in Section~\ref{sec:crossmachine} is
unaffected.} A model that is $82\%$ plant timetable still transfers between any
two machines sharing a shift pattern and still transfers to nothing else,
which is exactly what the transfer matrices measure. What the attribution adds
is that the machine's own history is not negligible: the previous production
run contributes $13.3\%$ of the attribution against $2.3\%$ of the gain, so the
forecaster is not purely a clock.

The chance-constraint classifier keys on something different, and the
difference is what justifies fitting it separately rather than deriving it
from the point forecast. Elapsed idle time is its single largest input
($33.1\%$ of attribution, rank 1) against sixth and $4.3\%$ for the duration
regressor. The feasibility question --- \emph{will this episode last at least
$h_{\min}$ longer?} --- depends on how long the episode has already run in a
way the conditional mean does not.

\begin{figure}[t]
  \centering
  \includegraphics[width=\linewidth]{fig_task13a_decision_shap_phase4}
  \caption{Shapley attribution of the remaining-idle-duration forecaster on
  the minimum-dwell labelling. (A)~Per-prediction attributions on $7\,440$
  held-out decision epochs; the horizontal axis is symmetric-log because
  attributions span $\pm10^{5}$~s while the bulk lies inside $\pm10^{3}$~s.
  Colour is the feature's own value. (B)~Mean $|\phi|$ against split gain; the
  dashed line is where the two measures would agree.}
  \label{fig:shap-decision}
\end{figure}

\subsection{The window forecaster is an electrical model, and mostly a
persistence one}
\label{sec:explain:window}

The window forecaster and the duration forecaster are separate models serving
separate consumers, and their attributions have almost nothing in common. The
window model draws $40.8\%$ of its attribution from the raw electrical
channels, $23.9\%$ from the rolling statistics and $20.1\%$ from the derived
electrical features, against $11.6\%$ for everything calendar
(Figure~\ref{fig:shap-window}A). Set beside the $82\%$ of the previous
subsection, this is the quantitative form of the observation in
Section~\ref{sec:crossmachine} that the two layers depend on nearly disjoint
information --- and it explains why one transfers between machines and the
other does not.

Two results here are worth stating plainly because neither is convenient.

\textbf{Voltage is the single largest contributor}, at $20.0\%$ of the
attribution across its nine design columns and $11.8\%$ on the final observed
value alone, against $3.0\%$ of the split gain. Supply voltage is not a state
variable of the machine; it sags under load and recovers when the drive
unloads, so it is a legitimate but indirect indicator, and one that would not
survive a change of supply transformer or a different point of common
coupling. A deployment on another site should expect this contribution not to
transfer.

\textbf{Almost half the attribution sits on the final observed value.}
Grouping the design columns by what they summarise
(Figure~\ref{fig:shap-window}B) gives $43.5\%$ to the last sample of the
lookback window, $26.4\%$ to block means and $26.3\%$ to block standard
deviations. A model whose largest single input is the state the machine is in
right now is a refinement of ``nothing changes'', and it cannot be expected to
beat that reference by much. It does not: the measured margin over untrained
persistence is $1.45$~s of per-window MAE
(Section~\ref{sec:significance}), which the attribution here now explains
rather than merely reports.

\begin{figure}[t]
  \centering
  \includegraphics[width=\linewidth]{fig_task13a_window_shap}
  \caption{Shapley attribution of the window state forecaster, refitted at
  seed~0 under the protocol of Section~\ref{sec:forecasting} and verified to
  reproduce its stored metrics exactly. (A)~The twelve largest measured
  channels, summed over their nine design columns each; hollow bars are the
  split gain. (B)~The same attribution grouped by what each design column
  summarises.}
  \label{fig:shap-window}
\end{figure}

\subsection{The state layer separates STANDBY on power factor}
\label{sec:explain:state}

Methods~1 and 2 of the five-method proof rest on a physical assertion: an
energised but unloaded induction motor is identified by its power factor and
its magnetising-current ratio, not by its power magnitude. That assertion is
testable against the labels the pipeline actually produced, by fitting a
gradient-boosted surrogate on them and attributing it.

On the hard split --- STANDBY against $\{$WORKING, PEAK\_LOAD$\}$, restricted
to energised rows so that the trivial OFF boundary cannot dominate --- the
attribution ranks \texttt{current} first ($17.8\%$), \texttt{power\_factor}
second ($16.9\%$) and \texttt{active\_power} third ($12.3\%$). The four
physics-named channels carry $30.9\%$ of the attribution between them.

The result is clearest where the two measures disagree. By split gain,
\texttt{active\_power} takes $60.3\%$ and \texttt{power\_factor} $5.1\%$; by
Shapley attribution the order reverses. \textbf{Power magnitude is what the
trees split on; power factor is what changes the predictions.} That is the
physics argument confirmed by the data rather than asserted from motor theory
--- and it is confirmed on the labelling the paper now quotes, not only on the
one the argument was originally made against. Agreement with the corresponding
Phase~II ranking on the earlier labelling is Spearman $\rho=0.833$
($p=4\times10^{-7}$) with four of five top features shared, so the
minimum-dwell re-labelling does not move the explanation.

Figure~\ref{fig:shap-state}A shows the mechanism directly: for
\texttt{power\_factor} and \texttt{current}, the attributions pushing towards
STANDBY are the low feature values, which is the magnetising-current signature
stated in Section~\ref{sec:method:accept_reject}, read off the model.

The surrogate reproduces its targets to $99.74\%$ accuracy. \emph{That number
is not a validation result and should not be read as one}: the targets are the
state model's own labels, so the accuracy measures only that a tree can
re-express a mixture model's decision boundary. Only the ranking is a result.

\begin{figure}[t]
  \centering
  \includegraphics[width=\linewidth]{fig_task13b_state_shap_phase4}
  \caption{Attribution of a gradient-boosted surrogate fitted on the
  GMM--HMM labels, for the STANDBY-versus-productive split on energised rows.
  (A)~Per-row attributions; blue is a low feature value.
  (B)~Three importance measures with independent failure modes, each
  normalised to its own total.}
  \label{fig:shap-state}
\end{figure}

\subsection{The transition structure}
\label{sec:explain:transitions}

The temporal model is claimed to learn a physically admissible state graph.
Figure~\ref{fig:transitions} reports three matrices that are commonly
conflated: the HMM's estimated transition matrix, the frequencies realised by
Viterbi decoding, and the frequencies after the minimum-dwell constraint.
Transitions are never counted across an acquisition gap.

\textbf{The prediction the state definition makes is that a machine cannot
move between de-energised and loaded in one second.} On pelletizer~I the
learned probability of OFF\,$\rightarrow$\,PEAK\_LOAD is exactly zero, and the
decoded sequence contains \emph{zero} such transitions in $5.47$~million
readings. OFF\,$\rightarrow$\,WORKING occurs $126$ times
($6.1\times10^{-5}$ of the outgoing mass from OFF), which at one-second
sampling is the expected consequence of a machine crossing the
STANDBY band faster than the sampling interval rather than a
violation of the ordering.

\textbf{The minimum-dwell constraint acts on the transitions, not only on the
labels.} WORKING\,$\rightarrow$\,OFF falls from $124$ events to $25$, and the
expected dwell in STANDBY rises from $96$~s to $246$~s. The constraint is
therefore doing what it was introduced to do --- removing runs too short to be
physical --- rather than merely relabelling rows.

It also introduces a small artefact that should be stated: three
PEAK\_LOAD\,$\rightarrow$\,OFF transitions appear after the constraint that
the decoder never produced. Absorbing a sub-threshold intermediate run into
its neighbours can join two states that were not adjacent in the decoded
sequence. Three events in $5.47$~million is $5\times10^{-7}$ of the record and
changes no reported quantity, but it is a property of the operator rather than
of the machine.

Across the plant the structure holds where the state model does and fails where
it does not. The largest forbidden-transition probability among the four
motor-driven machines is $1.9\times10^{-4}$ (pelletizer~I's
WORKING\,$\rightarrow$\,OFF); on exhaust-fan~I it is $3.4\times10^{-3}$,
eighteen times larger --- the same machine whose state names
Section~\ref{sec:multimachine} showed to be mis-assigned, arriving here through
an independent measurement.

\begin{figure}[t]
  \centering
  \includegraphics[width=\linewidth]{fig_task13c_transition_matrix}
  \caption{Transition structure on pelletizer~I. (A)~The HMM's estimated
  matrix. (B)~Frequencies counted from the Viterbi-decoded sequence.
  (C)~The same after the minimum-dwell constraint. Colour is logarithmic;
  parenthesised numbers are event counts for the transitions the state
  definition rules out.}
  \label{fig:transitions}
\end{figure}

%  Figures are referenced by bare name; add to the preamble
%      \graphicspath{{../outputs/phase6/figures/}}
%  and use the PDF copies (both PDF and PNG are written).
%  Bibliography: paper/references.bib
%  Numbers below are generated by  python run_phase6.py  and stored in
%    outputs/phase6/task13a/<machine>/task13a_results.json
%    outputs/phase6/task13b/<machine>/task13b_results.json
%    outputs/phase6/task13c/task13c_results.json
%  Figures:
%    fig_task13a_decision_shap_phase4, fig_task13a_window_shap,
%    fig_task13b_state_shap_phase4, fig_task13c_transition_matrix,
%    fig_task13c_transition_graph
%
%  FORWARD REFERENCES.  This file cites four sections that do not yet
%  have a .tex of their own; they will resolve to ?? until those are
%  written, which is deliberate rather than an oversight.  The labels
%  expected are:
%    sec:forecasting    the Task 5 / Task 12A forecasting comparison
%    sec:significance   the Phase V significance testing
%    sec:multimachine   the Task 8 eight-machine analysis
%    sec:crossmachine   the Task 9 transfer study
%  Everything else resolves against method_gmm_labeller.tex,
%  related_work.tex and decision_layer.tex as they stand.
% =====================================================================

\section{What the Models Use}
\label{sec:explain}

The three stages of the framework are opaque in different ways. The state
layer assigns a label without saying which measured quantity carried the
decision; both forecasters report split-gain importances, which are known to
be biased toward high-cardinality splits and describe the fitted trees rather
than any individual prediction. This section replaces those importances with
Shapley attributions~\cite{lundberg2017shap}, which are additive, consistent,
and defined per prediction, and reports the transition structure the temporal
model learns.

One caveat governs everything below and is not repeated at each result. Every
target attributed here is derived from the pipeline's own GMM--HMM labels.
These attributions therefore explain what each model learned; they are not
evidence that the underlying state assignment is correct. Validation of the
state definition is the physics battery of
Section~\ref{sec:method:accept_reject} and the cross-machine evidence of
Section~\ref{sec:crossmachine}, not this section.

\subsection{Attribution method}
\label{sec:explain:method}

Every model attributed here is a gradient-boosted
tree~\cite{chen2016xgboost} --- the two forecasters, the chance-constraint
classifier, and the surrogate fitted to the state labels --- so exact tree SHAP
applies and no sampling approximation is needed. Attributions are
computed on held-out rows: $7\,440$ decision epochs for the duration
forecaster, $15\,000$ design rows for the window forecaster, and $20\,000$
energised rows for the state surrogate.

Shapley values satisfy $f(x)=\mathbb{E}[f]+\sum_j\phi_j(x)$, and we verify that
identity rather than assume it. The verification is performed on the model's
raw margin --- the quantity tree SHAP decomposes --- and on a \emph{relative}
scale: the duration forecaster's outputs are seconds and run to $10^{5}$,
while the tree traversal accumulates in single precision, so an absolute
tolerance written for probability outputs is not the right instrument. The
largest relative reconstruction error measured over the four forecaster
attributions is $2.0\times10^{-6}$, four orders of magnitude below the
smallest difference any claim below rests on; the state surrogate's
attributions satisfy the stricter absolute check unmodified.

\subsection{The decision layer's forecaster is a calendar model}
\label{sec:explain:decision}

Section~\ref{sec:crossmachine} reported that the remaining-duration forecaster
draws $94.5\%$ of its split gain from four calendar features, and inferred
from that a model carrying the plant timetable rather than machine physics.
The Shapley attribution confirms the conclusion and corrects the number.

Under the minimum-dwell labelling the seven calendar features carry $82.4\%$
of the mean $|\phi|$ against $96.7\%$ of the split gain; on the
Phase~I--III labelling the same quantities are $78.0\%$ and $96.2\%$. The
direction is the same on both labellings, and the gap between the two measures
is systematic rather than incidental: gain and Shapley agree only at
Spearman $\rho=0.72$ over the ten features ($\rho=0.43$ on the Phase~I--III
labelling), and the two that move furthest are the two the original claim
leaned on.
Figure~\ref{fig:shap-decision}B makes the disagreement visible. The binary
\texttt{cur\_is\_weekend} takes $63.2\%$ of the split gain and $5.1\%$ of the
attribution: it is split on early and often, near the root, on a partition that
subsequent splits then refine, so it accumulates gain out of proportion to the
change it makes to any prediction. Conversely \texttt{onset\_hour} carries
$17.5\%$ of the gain and $43.1\%$ of the attribution.

\textbf{The corrected statement of the finding is therefore that four fifths,
not nineteen twentieths, of the decision forecaster's behaviour is calendar
--- and the conclusion drawn from it in Section~\ref{sec:crossmachine} is
unaffected.} A model that is $82\%$ plant timetable still transfers between any
two machines sharing a shift pattern and still transfers to nothing else,
which is exactly what the transfer matrices measure. What the attribution adds
is that the machine's own history is not negligible: the previous production
run contributes $13.3\%$ of the attribution against $2.3\%$ of the gain, so the
forecaster is not purely a clock.

The chance-constraint classifier keys on something different, and the
difference is what justifies fitting it separately rather than deriving it
from the point forecast. Elapsed idle time is its single largest input
($33.1\%$ of attribution, rank 1) against sixth and $4.3\%$ for the duration
regressor. The feasibility question --- \emph{will this episode last at least
$h_{\min}$ longer?} --- depends on how long the episode has already run in a
way the conditional mean does not.

\begin{figure}[t]
  \centering
  \includegraphics[width=\linewidth]{fig_task13a_decision_shap_phase4}
  \caption{Shapley attribution of the remaining-idle-duration forecaster on
  the minimum-dwell labelling. (A)~Per-prediction attributions on $7\,440$
  held-out decision epochs; the horizontal axis is symmetric-log because
  attributions span $\pm10^{5}$~s while the bulk lies inside $\pm10^{3}$~s.
  Colour is the feature's own value. (B)~Mean $|\phi|$ against split gain; the
  dashed line is where the two measures would agree.}
  \label{fig:shap-decision}
\end{figure}

\subsection{The window forecaster is an electrical model, and mostly a
persistence one}
\label{sec:explain:window}

The window forecaster and the duration forecaster are separate models serving
separate consumers, and their attributions have almost nothing in common. The
window model draws $40.8\%$ of its attribution from the raw electrical
channels, $23.9\%$ from the rolling statistics and $20.1\%$ from the derived
electrical features, against $11.6\%$ for everything calendar
(Figure~\ref{fig:shap-window}A). Set beside the $82\%$ of the previous
subsection, this is the quantitative form of the observation in
Section~\ref{sec:crossmachine} that the two layers depend on nearly disjoint
information --- and it explains why one transfers between machines and the
other does not.

Two results here are worth stating plainly because neither is convenient.

\textbf{Voltage is the single largest contributor}, at $20.0\%$ of the
attribution across its nine design columns and $11.8\%$ on the final observed
value alone, against $3.0\%$ of the split gain. Supply voltage is not a state
variable of the machine; it sags under load and recovers when the drive
unloads, so it is a legitimate but indirect indicator, and one that would not
survive a change of supply transformer or a different point of common
coupling. A deployment on another site should expect this contribution not to
transfer.

\textbf{Almost half the attribution sits on the final observed value.}
Grouping the design columns by what they summarise
(Figure~\ref{fig:shap-window}B) gives $43.5\%$ to the last sample of the
lookback window, $26.4\%$ to block means and $26.3\%$ to block standard
deviations. A model whose largest single input is the state the machine is in
right now is a refinement of ``nothing changes'', and it cannot be expected to
beat that reference by much. It does not: the measured margin over untrained
persistence is $1.45$~s of per-window MAE
(Section~\ref{sec:significance}), which the attribution here now explains
rather than merely reports.

\begin{figure}[t]
  \centering
  \includegraphics[width=\linewidth]{fig_task13a_window_shap}
  \caption{Shapley attribution of the window state forecaster, refitted at
  seed~0 under the protocol of Section~\ref{sec:forecasting} and verified to
  reproduce its stored metrics exactly. (A)~The twelve largest measured
  channels, summed over their nine design columns each; hollow bars are the
  split gain. (B)~The same attribution grouped by what each design column
  summarises.}
  \label{fig:shap-window}
\end{figure}

\subsection{The state layer separates STANDBY on power factor}
\label{sec:explain:state}

Methods~1 and 2 of the five-method proof rest on a physical assertion: an
energised but unloaded induction motor is identified by its power factor and
its magnetising-current ratio, not by its power magnitude. That assertion is
testable against the labels the pipeline actually produced, by fitting a
gradient-boosted surrogate on them and attributing it.

On the hard split --- STANDBY against $\{$WORKING, PEAK\_LOAD$\}$, restricted
to energised rows so that the trivial OFF boundary cannot dominate --- the
attribution ranks \texttt{current} first ($17.8\%$), \texttt{power\_factor}
second ($16.9\%$) and \texttt{active\_power} third ($12.3\%$). The four
physics-named channels carry $30.9\%$ of the attribution between them.

The result is clearest where the two measures disagree. By split gain,
\texttt{active\_power} takes $60.3\%$ and \texttt{power\_factor} $5.1\%$; by
Shapley attribution the order reverses. \textbf{Power magnitude is what the
trees split on; power factor is what changes the predictions.} That is the
physics argument confirmed by the data rather than asserted from motor theory
--- and it is confirmed on the labelling the paper now quotes, not only on the
one the argument was originally made against. Agreement with the corresponding
Phase~II ranking on the earlier labelling is Spearman $\rho=0.833$
($p=4\times10^{-7}$) with four of five top features shared, so the
minimum-dwell re-labelling does not move the explanation.

Figure~\ref{fig:shap-state}A shows the mechanism directly: for
\texttt{power\_factor} and \texttt{current}, the attributions pushing towards
STANDBY are the low feature values, which is the magnetising-current signature
stated in Section~\ref{sec:method:accept_reject}, read off the model.

The surrogate reproduces its targets to $99.74\%$ accuracy. \emph{That number
is not a validation result and should not be read as one}: the targets are the
state model's own labels, so the accuracy measures only that a tree can
re-express a mixture model's decision boundary. Only the ranking is a result.

\begin{figure}[t]
  \centering
  \includegraphics[width=\linewidth]{fig_task13b_state_shap_phase4}
  \caption{Attribution of a gradient-boosted surrogate fitted on the
  GMM--HMM labels, for the STANDBY-versus-productive split on energised rows.
  (A)~Per-row attributions; blue is a low feature value.
  (B)~Three importance measures with independent failure modes, each
  normalised to its own total.}
  \label{fig:shap-state}
\end{figure}

\subsection{The transition structure}
\label{sec:explain:transitions}

The temporal model is claimed to learn a physically admissible state graph.
Figure~\ref{fig:transitions} reports three matrices that are commonly
conflated: the HMM's estimated transition matrix, the frequencies realised by
Viterbi decoding, and the frequencies after the minimum-dwell constraint.
Transitions are never counted across an acquisition gap.

\textbf{The prediction the state definition makes is that a machine cannot
move between de-energised and loaded in one second.} On pelletizer~I the
learned probability of OFF\,$\rightarrow$\,PEAK\_LOAD is exactly zero, and the
decoded sequence contains \emph{zero} such transitions in $5.47$~million
readings. OFF\,$\rightarrow$\,WORKING occurs $126$ times
($6.1\times10^{-5}$ of the outgoing mass from OFF), which at one-second
sampling is the expected consequence of a machine crossing the
STANDBY band faster than the sampling interval rather than a
violation of the ordering.

\textbf{The minimum-dwell constraint acts on the transitions, not only on the
labels.} WORKING\,$\rightarrow$\,OFF falls from $124$ events to $25$, and the
expected dwell in STANDBY rises from $96$~s to $246$~s. The constraint is
therefore doing what it was introduced to do --- removing runs too short to be
physical --- rather than merely relabelling rows.

It also introduces a small artefact that should be stated: three
PEAK\_LOAD\,$\rightarrow$\,OFF transitions appear after the constraint that
the decoder never produced. Absorbing a sub-threshold intermediate run into
its neighbours can join two states that were not adjacent in the decoded
sequence. Three events in $5.47$~million is $5\times10^{-7}$ of the record and
changes no reported quantity, but it is a property of the operator rather than
of the machine.

Across the plant the structure holds where the state model does and fails where
it does not. The largest forbidden-transition probability among the four
motor-driven machines is $1.9\times10^{-4}$ (pelletizer~I's
WORKING\,$\rightarrow$\,OFF); on exhaust-fan~I it is $3.4\times10^{-3}$,
eighteen times larger --- the same machine whose state names
Section~\ref{sec:multimachine} showed to be mis-assigned, arriving here through
an independent measurement.

\begin{figure}[t]
  \centering
  \includegraphics[width=\linewidth]{fig_task13c_transition_matrix}
  \caption{Transition structure on pelletizer~I. (A)~The HMM's estimated
  matrix. (B)~Frequencies counted from the Viterbi-decoded sequence.
  (C)~The same after the minimum-dwell constraint. Colour is logarithmic;
  parenthesised numbers are event counts for the transitions the state
  definition rules out.}
  \label{fig:transitions}
\end{figure}

%  Figures are referenced by bare name; add to the preamble
%      \graphicspath{{../outputs/phase6/figures/}}
%  and use the PDF copies (both PDF and PNG are written).
%  Bibliography: paper/references.bib
%  Numbers below are generated by  python run_phase6.py  and stored in
%    outputs/phase6/task13a/<machine>/task13a_results.json
%    outputs/phase6/task13b/<machine>/task13b_results.json
%    outputs/phase6/task13c/task13c_results.json
%  Figures:
%    fig_task13a_decision_shap_phase4, fig_task13a_window_shap,
%    fig_task13b_state_shap_phase4, fig_task13c_transition_matrix,
%    fig_task13c_transition_graph
%
%  FORWARD REFERENCES.  This file cites four sections that do not yet
%  have a .tex of their own; they will resolve to ?? until those are
%  written, which is deliberate rather than an oversight.  The labels
%  expected are:
%    sec:forecasting    the Task 5 / Task 12A forecasting comparison
%    sec:significance   the Phase V significance testing
%    sec:multimachine   the Task 8 eight-machine analysis
%    sec:crossmachine   the Task 9 transfer study
%  Everything else resolves against method_gmm_labeller.tex,
%  related_work.tex and decision_layer.tex as they stand.
% =====================================================================

\section{What the Models Use}
\label{sec:explain}

The three stages of the framework are opaque in different ways. The state
layer assigns a label without saying which measured quantity carried the
decision; both forecasters report split-gain importances, which are known to
be biased toward high-cardinality splits and describe the fitted trees rather
than any individual prediction. This section replaces those importances with
Shapley attributions~\cite{lundberg2017shap}, which are additive, consistent,
and defined per prediction, and reports the transition structure the temporal
model learns.

One caveat governs everything below and is not repeated at each result. Every
target attributed here is derived from the pipeline's own GMM--HMM labels.
These attributions therefore explain what each model learned; they are not
evidence that the underlying state assignment is correct. Validation of the
state definition is the physics battery of
Section~\ref{sec:method:accept_reject} and the cross-machine evidence of
Section~\ref{sec:crossmachine}, not this section.

\subsection{Attribution method}
\label{sec:explain:method}

Every model attributed here is a gradient-boosted
tree~\cite{chen2016xgboost} --- the two forecasters, the chance-constraint
classifier, and the surrogate fitted to the state labels --- so exact tree SHAP
applies and no sampling approximation is needed. Attributions are
computed on held-out rows: $7\,440$ decision epochs for the duration
forecaster, $15\,000$ design rows for the window forecaster, and $20\,000$
energised rows for the state surrogate.

Shapley values satisfy $f(x)=\mathbb{E}[f]+\sum_j\phi_j(x)$, and we verify that
identity rather than assume it. The verification is performed on the model's
raw margin --- the quantity tree SHAP decomposes --- and on a \emph{relative}
scale: the duration forecaster's outputs are seconds and run to $10^{5}$,
while the tree traversal accumulates in single precision, so an absolute
tolerance written for probability outputs is not the right instrument. The
largest relative reconstruction error measured over the four forecaster
attributions is $2.0\times10^{-6}$, four orders of magnitude below the
smallest difference any claim below rests on; the state surrogate's
attributions satisfy the stricter absolute check unmodified.

\subsection{The decision layer's forecaster is a calendar model}
\label{sec:explain:decision}

Section~\ref{sec:crossmachine} reported that the remaining-duration forecaster
draws $94.5\%$ of its split gain from four calendar features, and inferred
from that a model carrying the plant timetable rather than machine physics.
The Shapley attribution confirms the conclusion and corrects the number.

Under the minimum-dwell labelling the seven calendar features carry $82.4\%$
of the mean $|\phi|$ against $96.7\%$ of the split gain; on the
Phase~I--III labelling the same quantities are $78.0\%$ and $96.2\%$. The
direction is the same on both labellings, and the gap between the two measures
is systematic rather than incidental: gain and Shapley agree only at
Spearman $\rho=0.72$ over the ten features ($\rho=0.43$ on the Phase~I--III
labelling), and the two that move furthest are the two the original claim
leaned on.
Figure~\ref{fig:shap-decision}B makes the disagreement visible. The binary
\texttt{cur\_is\_weekend} takes $63.2\%$ of the split gain and $5.1\%$ of the
attribution: it is split on early and often, near the root, on a partition that
subsequent splits then refine, so it accumulates gain out of proportion to the
change it makes to any prediction. Conversely \texttt{onset\_hour} carries
$17.5\%$ of the gain and $43.1\%$ of the attribution.

\textbf{The corrected statement of the finding is therefore that four fifths,
not nineteen twentieths, of the decision forecaster's behaviour is calendar
--- and the conclusion drawn from it in Section~\ref{sec:crossmachine} is
unaffected.} A model that is $82\%$ plant timetable still transfers between any
two machines sharing a shift pattern and still transfers to nothing else,
which is exactly what the transfer matrices measure. What the attribution adds
is that the machine's own history is not negligible: the previous production
run contributes $13.3\%$ of the attribution against $2.3\%$ of the gain, so the
forecaster is not purely a clock.

The chance-constraint classifier keys on something different, and the
difference is what justifies fitting it separately rather than deriving it
from the point forecast. Elapsed idle time is its single largest input
($33.1\%$ of attribution, rank 1) against sixth and $4.3\%$ for the duration
regressor. The feasibility question --- \emph{will this episode last at least
$h_{\min}$ longer?} --- depends on how long the episode has already run in a
way the conditional mean does not.

\begin{figure}[t]
  \centering
  \includegraphics[width=\linewidth]{fig_task13a_decision_shap_phase4}
  \caption{Shapley attribution of the remaining-idle-duration forecaster on
  the minimum-dwell labelling. (A)~Per-prediction attributions on $7\,440$
  held-out decision epochs; the horizontal axis is symmetric-log because
  attributions span $\pm10^{5}$~s while the bulk lies inside $\pm10^{3}$~s.
  Colour is the feature's own value. (B)~Mean $|\phi|$ against split gain; the
  dashed line is where the two measures would agree.}
  \label{fig:shap-decision}
\end{figure}

\subsection{The window forecaster is an electrical model, and mostly a
persistence one}
\label{sec:explain:window}

The window forecaster and the duration forecaster are separate models serving
separate consumers, and their attributions have almost nothing in common. The
window model draws $40.8\%$ of its attribution from the raw electrical
channels, $23.9\%$ from the rolling statistics and $20.1\%$ from the derived
electrical features, against $11.6\%$ for everything calendar
(Figure~\ref{fig:shap-window}A). Set beside the $82\%$ of the previous
subsection, this is the quantitative form of the observation in
Section~\ref{sec:crossmachine} that the two layers depend on nearly disjoint
information --- and it explains why one transfers between machines and the
other does not.

Two results here are worth stating plainly because neither is convenient.

\textbf{Voltage is the single largest contributor}, at $20.0\%$ of the
attribution across its nine design columns and $11.8\%$ on the final observed
value alone, against $3.0\%$ of the split gain. Supply voltage is not a state
variable of the machine; it sags under load and recovers when the drive
unloads, so it is a legitimate but indirect indicator, and one that would not
survive a change of supply transformer or a different point of common
coupling. A deployment on another site should expect this contribution not to
transfer.

\textbf{Almost half the attribution sits on the final observed value.}
Grouping the design columns by what they summarise
(Figure~\ref{fig:shap-window}B) gives $43.5\%$ to the last sample of the
lookback window, $26.4\%$ to block means and $26.3\%$ to block standard
deviations. A model whose largest single input is the state the machine is in
right now is a refinement of ``nothing changes'', and it cannot be expected to
beat that reference by much. It does not: the measured margin over untrained
persistence is $1.45$~s of per-window MAE
(Section~\ref{sec:significance}), which the attribution here now explains
rather than merely reports.

\begin{figure}[t]
  \centering
  \includegraphics[width=\linewidth]{fig_task13a_window_shap}
  \caption{Shapley attribution of the window state forecaster, refitted at
  seed~0 under the protocol of Section~\ref{sec:forecasting} and verified to
  reproduce its stored metrics exactly. (A)~The twelve largest measured
  channels, summed over their nine design columns each; hollow bars are the
  split gain. (B)~The same attribution grouped by what each design column
  summarises.}
  \label{fig:shap-window}
\end{figure}

\subsection{The state layer separates STANDBY on power factor}
\label{sec:explain:state}

Methods~1 and 2 of the five-method proof rest on a physical assertion: an
energised but unloaded induction motor is identified by its power factor and
its magnetising-current ratio, not by its power magnitude. That assertion is
testable against the labels the pipeline actually produced, by fitting a
gradient-boosted surrogate on them and attributing it.

On the hard split --- STANDBY against $\{$WORKING, PEAK\_LOAD$\}$, restricted
to energised rows so that the trivial OFF boundary cannot dominate --- the
attribution ranks \texttt{current} first ($17.8\%$), \texttt{power\_factor}
second ($16.9\%$) and \texttt{active\_power} third ($12.3\%$). The four
physics-named channels carry $30.9\%$ of the attribution between them.

The result is clearest where the two measures disagree. By split gain,
\texttt{active\_power} takes $60.3\%$ and \texttt{power\_factor} $5.1\%$; by
Shapley attribution the order reverses. \textbf{Power magnitude is what the
trees split on; power factor is what changes the predictions.} That is the
physics argument confirmed by the data rather than asserted from motor theory
--- and it is confirmed on the labelling the paper now quotes, not only on the
one the argument was originally made against. Agreement with the corresponding
Phase~II ranking on the earlier labelling is Spearman $\rho=0.833$
($p=4\times10^{-7}$) with four of five top features shared, so the
minimum-dwell re-labelling does not move the explanation.

Figure~\ref{fig:shap-state}A shows the mechanism directly: for
\texttt{power\_factor} and \texttt{current}, the attributions pushing towards
STANDBY are the low feature values, which is the magnetising-current signature
stated in Section~\ref{sec:method:accept_reject}, read off the model.

The surrogate reproduces its targets to $99.74\%$ accuracy. \emph{That number
is not a validation result and should not be read as one}: the targets are the
state model's own labels, so the accuracy measures only that a tree can
re-express a mixture model's decision boundary. Only the ranking is a result.

\begin{figure}[t]
  \centering
  \includegraphics[width=\linewidth]{fig_task13b_state_shap_phase4}
  \caption{Attribution of a gradient-boosted surrogate fitted on the
  GMM--HMM labels, for the STANDBY-versus-productive split on energised rows.
  (A)~Per-row attributions; blue is a low feature value.
  (B)~Three importance measures with independent failure modes, each
  normalised to its own total.}
  \label{fig:shap-state}
\end{figure}

\subsection{The transition structure}
\label{sec:explain:transitions}

The temporal model is claimed to learn a physically admissible state graph.
Figure~\ref{fig:transitions} reports three matrices that are commonly
conflated: the HMM's estimated transition matrix, the frequencies realised by
Viterbi decoding, and the frequencies after the minimum-dwell constraint.
Transitions are never counted across an acquisition gap.

\textbf{The prediction the state definition makes is that a machine cannot
move between de-energised and loaded in one second.} On pelletizer~I the
learned probability of OFF\,$\rightarrow$\,PEAK\_LOAD is exactly zero, and the
decoded sequence contains \emph{zero} such transitions in $5.47$~million
readings. OFF\,$\rightarrow$\,WORKING occurs $126$ times
($6.1\times10^{-5}$ of the outgoing mass from OFF), which at one-second
sampling is the expected consequence of a machine crossing the
STANDBY band faster than the sampling interval rather than a
violation of the ordering.

\textbf{The minimum-dwell constraint acts on the transitions, not only on the
labels.} WORKING\,$\rightarrow$\,OFF falls from $124$ events to $25$, and the
expected dwell in STANDBY rises from $96$~s to $246$~s. The constraint is
therefore doing what it was introduced to do --- removing runs too short to be
physical --- rather than merely relabelling rows.

It also introduces a small artefact that should be stated: three
PEAK\_LOAD\,$\rightarrow$\,OFF transitions appear after the constraint that
the decoder never produced. Absorbing a sub-threshold intermediate run into
its neighbours can join two states that were not adjacent in the decoded
sequence. Three events in $5.47$~million is $5\times10^{-7}$ of the record and
changes no reported quantity, but it is a property of the operator rather than
of the machine.

Across the plant the structure holds where the state model does and fails where
it does not. The largest forbidden-transition probability among the four
motor-driven machines is $1.9\times10^{-4}$ (pelletizer~I's
WORKING\,$\rightarrow$\,OFF); on exhaust-fan~I it is $3.4\times10^{-3}$,
eighteen times larger --- the same machine whose state names
Section~\ref{sec:multimachine} showed to be mis-assigned, arriving here through
an independent measurement.

\begin{figure}[t]
  \centering
  \includegraphics[width=\linewidth]{fig_task13c_transition_matrix}
  \caption{Transition structure on pelletizer~I. (A)~The HMM's estimated
  matrix. (B)~Frequencies counted from the Viterbi-decoded sequence.
  (C)~The same after the minimum-dwell constraint. Colour is logarithmic;
  parenthesised numbers are event counts for the transitions the state
  definition rules out.}
  \label{fig:transitions}
\end{figure}
